\chapter{Literature Review}
\label{chap:lr}
\chaptermark{Second Chapter Heading}

\section{Refinement Types}
\label{sec:refinement_types_defintion}
Refinement types allow the developer to extend a language's type system
with predicates \cite{jhalaRefinementTypesTutorial2020}. For example, the user can constraint
the regular integer type with the following refinement 
to get the type that represents only natural numbers.
\begin{lstlisting}
    @Refinement(identifier: 'x', predicate: 'x >= 0')
    typedef nat = int;
\end{lstlisting}
With such refinements the developer can easily avoid index out of bounds errors by, for example,
defining the following refinement:
\begin{lstlisting}
    @Refinement(
        parameters: ['list', 'index'],
        predicate: 'index >= 0 && index < length(list)',
    )
    int get(List<int> list, int index) => list[index];
\end{lstlisting}
The type checker will ensure that the index provided to the function is smaller than the
length of the list, and if not, will report this operation as unsafe.

Refinement types can also be used to express constraints to function parameters that 
cannot be captured by ordinary type checking. Consider a function that finds an average of
a list of integers. Although its parameter is correctly typed \texttt{List<int>}, the 
function cannot safely process an empty list, since it will involve division by zero:
\begin{lstlisting}
    @Refinement(
        identifier: 'values',
        predicate: 'length(values) > 0',
    )
    double calculateAverage(List<int> values) {
      return values.reduce((a, b) => a + b) / values.length;
    }

    calculateAverage([2, 4, 6]); // Safe
    calculateAverage([]);        // Refinement violation
\end{lstlisting}


\section{Liquid Types and Refinement Type Inference}
\label{sec:liquid_types}
Annotating each variable, each type with refinement information can get
overwhelming and consume more time than the annotations will reduce in error finding.
To address this issue we can implement refinement type inference using \textit{Logically Qualified Data Types} also
known as \textit{Liquid Types} system explored by \textit{Rondon et al.}\cite{rondonLiquidTypes2008}
which will provide the developer with strength of refinement types without additional unnecessary effort.
The inference process involves the following steps:
\begin{enumerate}
    \item Hindley-Milner type inference \cite{damasPrincipalTypeSchemesFunctional}
    \item Constraint generation
    \item Constraint solving using an SMT solver \cite{demouraZ3EfficientSMT2008}
\end{enumerate}
Assume a function \textit{safeHead}, where we want to check safety of \textit{xs[0]} operation:
\begin{lstlisting}
    int safeHead(List<int> xs) => xs.length > 0 ? xs[0] : 0;
\end{lstlisting}
Without type inference we would need to manually specify length property of \textit{xs}
as a refinement, however liquid types can infer the constraint from the program itself. 

\section{Refinement Types in Other Languages}
\label{sec:existing_implementations}
Similar implementations exist for other programming languages. The most similar one to my implementation is
Refined TypeScript (RSC) \cite{vekrisRefinementTypesTypeScript2016} by \textit{Vekris et al}.
Their work demonstrates an integration of a refinement type system into an imperative object-oriented language while
retaining support for higher-order functions.
\\
One of the main challenges addressed by RSC is interaction between refinements and mutable
program state. Since the value of a variable can change throughout runtime, the previously
established refinement can become invalid after reassignment. To address this issue the 
authors translate the input program into Static Single Assessment (SSA) form, where each variable
is only assigned once, which allows the checker to reason only about a single variable version
in context in which it was assigned.
\\
However, SSA alone does not solve the problem of mutation through references.
RSC addresses this issue by separating mutable and immutable fields.
Refinements are allowed to depend on immutable fields, since their value doesn't change
after initialization. Mutable fields are subject to class invariants that must be preserved
whenever they are updated.
\\
RSC demonstrates this approach using a two-dimensional array.
The dimension of the 2D array are speicified through refinement
predicates that relate its length to the immutable height and width of the containing object.
When the array value is replaced, the checker verifies that the new value satisfies the speicified
size constraints. This approach allows RSC to
verify array access safety and object invariants in cases of mutations. \cite{vekrisRefinementTypesTypeScript2016}

Another implementation that addresses mutable data types is Flux: Liquid Types for Rust \cite{lehmannFluxLiquidTypes2023}
by \textit{Lehmann et al.}.
Unlike RSC, Flux builds upon Rust's ownership and borrowing mechanisms to
analyze mutable memory locations. Its main idea is to combine
refinement types with ownership guarantees, allowing the checker to track
changes to the properties of mutable data without losing soundness due to
aliasing.

Flux distinguishes between exclusive ownership and borrowing. When a
function has exclusive ownership of a value, it can modify that value and
update its refinement accordingly. This is referred to as a strong update.
For example, a variable representing an integer with a refinement stating
that its value equals \(n\) can update this refinement to \(n+1\) after an
incrementation of the variable.

Borrowing introduces additional restrictions. A shared reference provides
read-only access, while a mutable reference permits modifications but
requires the original refinement invariant to remain satisfied. Flux also
introduces strong references, which allow functions to specify how the
refinement of the referenced value changes after a call.

Particularly relevant to this thesis is Flux's treatment of mutable
collections. The authors demonstrate a refined API for Rust vectors
(\texttt{RVec}), where the vector's length is represented directly in its
type. For example, a vector of length \(n\) has a type indexed by \(n\),
and its indexing operations require an index smaller than that
length. Operations that modify the vector can update this refinement.
For example the \texttt{push} operation, updates the vector length refinement to be \(n + 1\)
The following indexing operations are using this updated length information.

This approach is particularly relevant to the proposed implementation,
as it demonstrates how refinement types can be tracked across
operations that change list size. However, Flux relies on Rust's
ownership and borrowing guarantees to control aliasing and mutation.

Other notable implementations include LiquidHaskell
\cite{vazouLiquidHaskellExperienceRefinement2014} and Refinement Types
for Ruby \cite{kazerounianRefinementTypesRuby2017}. LiquidHaskell
demonstrates the application of refinement types to a functional language,
where immutable values simplify the preservation of
refinement predicates. In contrast, the Ruby implementation (RTR) addresses
mutable object fields by explicitly tracking which fields a function may
modify. When a method is treated modularly, the verifier invalidates
the previous knowledge about the modified fields and relies on the
method's postconditions to establish their properties after the call.
This is implemented through a \texttt{modifies} specification and
symbolic state updates.

These implementations demonstrate different approaches to handling
of mutable data, providing a useful foundation for the proposed implementation.


\section{SMT Solvers and Z3}
\label{sec:smt_solvers}
Satisfiability Modulo Theories (SMT) is a decision problem that asks whether a first-order logic
formula is satisfiable with respect to a background theory \cite{barrettChapter33Satisfiability2021}. It extends Boolean Satisfiability (SAT) problem
by adding reasoning such as integers, arrays, bit-verctos and uninterpreted functions. Since SMT solvers
are able of reasoning about relationships between values and operations within theories, they become
particularly useful in software verification, where program properties involve arithmetic constraints, equality connections
and data structures.

In context of refinement type checking, an SMT solver serves as an engine responsible for verifying
refinement predicates. The checker converts program properties into logical constraints, which are then
passed to the solver. For example to determine the safety of a list indexing operation, the implementation must determine
that the index is smaller than the list length and is larger or equal to zero. This can be expressed 
as a logical implication and determined truthful by checking whether the opposite is satisfiable.
Unsatisfied results means that the original property holds under given conditions.

Several SMT solvers can be used for this problem. In this section we are considering three popular
implementations: Z3, cvc5 and Yices 2. Although all three support required functionality, they differ in
their architecture, supported interfaces and additional capabilities.

\subsection{Z3}
Z3 is an SMT solver developed by Microsoft and introduced by de Moura and Bjørner \cite{demouraZ3EfficientSMT2008}.
It was designed with application in software verification and formal methods in mind. Its general-purpose
SMT engine follows the Conflict-Driven Clause Learning modulo (CDCL(T)) architecture. In this approach,
an SAT solver coordinates specialized theory solvers that exchange information between the constraints they are
solving.

Z3 supports a wide range of theories, including integer and real arithmetics, arrays, bit-vectors,
uninterpreted functions and algebraic datatypes. It lso provides APIs for several programming languages
including C, C++, Java, Python and .NET, as well as supports SMT-LIB input format. This makes it possible
to incorporate Z3 into any programming language by invoking an independent process communicationg through
SMT-LIB scripts.

Its popularity in software verification and extensive theory support make Z3 a suitable candidate for
the proposed refinement type checker.

\subsection{cvc5}
cvc5 is an open-source SMT solver that was developed as a versatile and extensible theorem-proving
engine, with particular focus on supporting a broad range of theories and proof production \cite{barbosaCvc5VersatileIndustrialStrength2022}.
Similar to Z3, its main solving architecture is based on CDCL(T).

cvc5 supports integer and real arithmetic, arrays, bit-vectors, algebraic datatypes, strings, sets and several
other theories. One of its unique features is its proof-production infrastructure, that is capable of generating
detailed proofs of unsatisfiability. This is particularly relevant in applications that require independently 
checkable verification results

The solver provides an API for C++, with additional interfaces for C, Java and Python, as well as can be
used with SMT-LIB. However, proof production is not a primary requirement of the proposed implementation,
which only needs to determine whether generated verification conditions are satisfiable.

\subsection{Yices 2}
Yices 2 is an SMT solver developed by SRI International and presented by Dutertre \cite{dutertreYices222014}.
Its architecture is modular and allows user to configure solving contexts with different solvers and solving procedures.
It supports the standard CDCL(T) approach for many theories as well as Model-Constructing Satisfiability Calculus
(MCSAT) solver for problems involving nonliear arithmetics.

Yices 2 supports integer and real arithmetic, arrays, bit-vectors, uninterpreted functions, and additional data types.
It accepts formulas in SMT-LIB format and provides a C API and bindings for Java, Python, Go, and OCaml.

An important characteristic of Yices 2 is its focus on efficient solving of arithmetic problems. These characteristics
make it a viable alternative for refinement type checking, particularly when verification conditions primarly
consist of arithmetic constraints.

\subsection{Comparison and Selection}
All three solvers provide the basic functionality required by the proposed refinement type checker.
Integer arithmetic and Boolean reasoning are sufficient for the initial verification scope.
This makes the choice of solver based mainly on integration capabilities, extensibility, and suitability
for software verification.

Z3 offers a convenient combination of these properties. Its established role in software verification, ease of implementation
via SMT-LIB scripts make it a practical choice for integration into the proposed implementation. Its architecture 
also provides specialized solving procedures that may become useful if the verification language 
scope is extended.

cvc5 offers similar general-purpose SMT capabilities and has strong proof-production infrastructure.
However, generating independently checkable proofs is not required by this project, making this advantage
not particularly important for the implementation. Yices 2 provides and efficient and configurable alternative,
especially for arithmetic-heavy problems, but its advantages are less significant for the limited set of theories
currently required.

Based on these conclusions, Z3 is selected as the SMT solver for the proposed refinement type checker.
This is primarly motivated by its focus on software verification, broad integration support and possibility
for painless future extensions.


% de Moura and Bjørner present Z3 \cite{demouraZ3EfficientSMT2008}, an SMT solver developed at Microsoft
% Research and designed specifically for applications in software verification. In context of refinement
% type checking, an SMT solver can serve as the reasoning engine, responsible for verifying refinement predicates and
% generating constraints, making Z3 suitable for proposed refinement type checker for Dart. Using Z3 allows to
% delegate logical reasoning to Z3 and focus on the implementation on extracting extracting properties and interpreting
% verification results.

% \Blindtext[2]

% \section{Another Section}
% \Blindtext[1]
% \newpage
% With the widespread of computing systems, information processing, and net working, the practice of replacing paper documentation to electronic documentation has become more and more common. Electronic documentation within the workplace has several advantages over the traditional one, such as being easy to share, copy and edit the document. However, these advantages also present a problem. A malefactor, having access to the system, can easily copy and leak the document, without leaving any trace. Such actions are virtually undetectable in most systems, so, the malefactor goes unpunished. In this thesis, we propose one solution to the problem: digital watermarking. 

% Every electronic document within the protected system is marked with an invisible digital watermark, containing information about the user, accessing this particular document. Therefore, in case the protected company discovers the leaked document, they will be able to identify the machine of the malicious person and time when the document was leaked. This will allow inflicting punishment on the malefactor, recovering the costs of the leak, and potentially preventing future ones. 

% \begin{longtable}{c|c}
% \caption[This is the title I want to appear in the List of Tables]{Simulation Parameters} \label{table:secsimulation_params} \\
% \hline
% A & B  \\
% \hline
% \endfirsthead
% \multicolumn{2}{@{}l}{} \\
% \hline
% A & B \\
% \hline
% \endhead
% \hline
%  \textbf{Parameter} & \textbf{Value}\\
%  \hline
%  Number of vehicles & $|\mathcal{V}|$\\
%  \hline
%  Number of RSUs & $|\mathcal{U}|$\\
%  \hline
%  RSU coverage radius & 150 m\\
%  \hline
%  V2V communication radius & 30 m\\
%  \hline
%  Smart vehicle antenna height & 1.5 m\\
%  \hline
%  RSU antenna height & 25 m\\
%  \hline
%  Smart vehicle maximum speed & $v_{max}$ m/s\\
%  \hline
%  Smart vehicle minimum speed & $v_{min}$ m/s\\
%  \hline
%  Common smart vehicle cache capacities & $[50, 100, 150, 200, 250]$ mb\\
%  \hline
%  Common RSU cache capacities & $[5000,1000,1500,2000,2500]$ mb\\
%  \hline
%  Common backhaul rates & $[75, 100, 150]$ mb/s\\
%  \hline
% \end{longtable}

% \begin{figure}[hbt]
% \centering
% \includegraphics[]{figs/images.png}
% \caption{One kernel at $x_s$ (\emph{dotted kernel}) or two kernels at
% $x_i$ and $x_j$ (\textit{left and right}) lead to the same summed estimate
% at $x_s$. This shows a figure consisting of different types of
% lines. Elements of the figure described in the caption should be set in
% italics, in parentheses, as shown in this sample caption.}
% \label{fig:secex}
% \end{figure}

% This description implies several essential properties of the task at hand:
% \begin{enumerate}
%     \item Watermark must contain all necessary information, but still, be placeable and recognizable even on smaller images. The produced watermark must be compact but have the possibility to store enough information. 
%     \item To prevent easy tampering, the watermark must be invisible to the naked eye (and, preferably, to basic image parsing tools). If malefactor does not know about the existence of watermark, they might not even try to remove it and disable it. 
% \end{enumerate}
